\documentclass[11pt,a4paper]{article}
\usepackage[utf8]{inputenc}\usepackage[T1]{fontenc}
\usepackage{lmodern,amsmath,amssymb,booktabs,geometry}
\usepackage[hidelinks]{hyperref}\geometry{margin=2.3cm}
\setlength{\parindent}{0pt}\setlength{\parskip}{0.5em}
\title{Radau: Numerical Method and Adaptive Execution}
\author{}\date{}\begin{document}\maketitle
\section{Differential equation and numerical method}
Consider $z'=f(t,z)$, $z(t_0)=z_0$, in the project's packed physical layout.
A three-stage Radau IIA collocation method of order five, with an embedded third-order error estimator and a cubic collocation interpolant. The implementation delegates the numerical tableau, local error
estimator, trial rejection and interpolation polynomial to the installed SciPy
\texttt{Radau} solver. The adapter does not duplicate these algorithms.

A Runge--Kutta evaluation has stage equations
\[
Y_i=z_n+h\sum_j a_{ij}f(t_n+c_jh,Y_j),\qquad
z_{n+1}=z_n+h\sum_i b_i f(t_n+c_ih,Y_i).
\]
The collocation equations are implicit and are solved by SciPy Newton iterations. The public adapter accepts an optional full-state Jacobian callable; otherwise SciPy uses finite differences. With energy tracking, a supplied Jacobian must cover all augmented coordinates.
\section{Error control and accepted intervals}
The adaptive backend scales local error estimates using positive absolute and
relative tolerances, with component scale based on
$\mathrm{atol}+\mathrm{rtol}\max(|z_n|,|z_{n+1}|)$. It selects subsequent step
sizes and can reject trial steps. \texttt{max\_step} bounds accepted duration;
\texttt{first\_step=None} delegates initial selection to the backend.
Requested tolerances bound neither global trajectory error nor roundoff.

One live solver is retained over the complete interval. Its state, stage work
and any Jacobian/factorization reuse survive accepted steps. The adapter calls
\texttt{solver.step()} until acceptance or failure; rejected trials never
become observer records or separate accepted metric rows.
\section{Dense output and energy extension}
Saved output times do not determine the accepted grid. Each requested sample
is evaluated from the dense interpolant of its accepted interval. This does not
restart the solver. \texttt{dense\_output=True} explicitly constructs every
accepted interpolant; optional observers can retain them for later time queries.

For a suitable time-dependent Hamiltonian system, optional tracking augments
$z$ with one momentum per particle:
\[
k'=-\partial_t H(t,z),\qquad k(t_0)=0.
\]
The exported generalized-energy diagnostic measures $H(t,z)+k$. Error control
acts on the augmented vector, so this option can change the accepted time grid
and the physical floating-point trajectory. It is not an independent proof of
trajectory accuracy.
\section{Shared architecture and observation}
The public class owns initialization, advance, observation and history export.
Inherited \texttt{new\_run} constructs a fresh instance of the same class and
calls \texttt{initialize}, which creates its live solver. There is no separate
method context or record of callbacks. All 13 methods inherit the same
integration entry point and use the common global coordinator.

\texttt{ScipyAdaptiveController} exposes accepted intervals from the retained
numerical session. The common collector stores requested samples and small
accepted-step statistics; it does not retain interpolants or events by default.

An \texttt{AdaptiveIntegrationStep} records times and copied states. It also
exposes work counters and \texttt{dense\_state}. It does not claim a fixed-duration
candidate-state map: adaptive decisions depend on numerical history.
\texttt{AdaptiveTrajectoryObserver} explicitly retains continuous output outside
the integrator. The observer and any persistence belong to the caller.

Function evaluations, Jacobian evaluations and LU decompositions are counted
per accepted interval. Initialization work belongs to the first row; rejected
trial work belongs to the acceptance that follows it. Interpolation needed only
for an observer is excluded from main counters unless continuous output was
explicitly requested. The ordinary method operation \texttt{advance} accepts one
step on the retained solver, taking an upper step bound. The controller reports
the actual accepted endpoint. A run instance belongs to one execution; repeating
\texttt{simulate} creates a fresh instance and solver.
\section{Implementation and verification}
The adapter is in \path{src/simulation/methods/adaptive/scipy.py}; shared
coordination and collection are in \path{src/simulation/integration.py}.
The model's six-phase diagram is documented in its \texttt{simulation/}
directory. Tests compare sampled states and work counters with the same
SciPy solver, check a known non-autonomous rotation, and verify optional energy,
dense retention, repeated runs, observation isolation and failure handling.
Scientific reference trajectories still require refinement and an independent
audit; arithmetic precision and integration error remain separate.
\end{document}
